% ============================================================================
% Fichier : annexes/Ann_A_glossaire.tex
% Titre   : Glossaire
% Version : 1.0 -- novembre 2026
% ============================================================================
% Définitions reprises, en abrégé, des encadrés « Définition » des chapitres
% (termes marqués \termecle) ; renvoi au chapitre de référence.
% ============================================================================

\chapter{Glossaire}
\label{ann:glossaire}

Ce glossaire reprend, sous une forme abrégée, les définitions données dans
les chapitres. Le renvoi indique où le terme est défini et illustré.

\begin{description}
    \item[Accord de niveau de service (SLA)] Accord documenté entre un
          fournisseur et son client, qui précise le service et les niveaux
          attendus (Ch.~\ref{chap:itil-niveaux}).
    \item[Accord de niveau opérationnel (OLA)] Accord entre le fournisseur et
          une autre équipe interne qui soutient le service
          (Ch.~\ref{chap:itil-niveaux}).
    \item[Activité (COBIT)] Action concrète qui met en œuvre une pratique~;
          rattachée à un niveau de capacité de~2 à~5
          (Ch.~\ref{chap:cobit-objectifs}).
    \item[Alignement stratégique] Degré de cohérence entre la stratégie de
          l'organisation, celle des technologies de l'information, et les
          infrastructures et processus qui les mettent en œuvre
          (Ch.~\ref{chap:alignement}).
    \item[Alignement volontaire] Choix de s'inspirer d'un référentiel ou
          d'une norme sans viser de certificat (Ch.~\ref{chap:normes}).
    \item[Analyse d'impact sur l'activité (BIA)] Analyse qui identifie les
          activités essentielles, mesure l'effet de leur interruption dans
          le temps et en déduit les objectifs de reprise
          (Ch.~\ref{chap:itil-niveaux}).
    \item[Autorité de changement] Personne ou groupe qui autorise un
          changement~; l'autorisation est déléguée selon le risque
          (Ch.~\ref{chap:itil-changements}).
    \item[Base de données de gestion de la configuration (CMDB)] Base qui
          enregistre les éléments de configuration, leurs attributs et
          leurs relations (Ch.~\ref{chap:itil-changements}).
    \item[Cadre de référence] Ensemble structuré de principes et
          d'objectifs qui guide un domaine sans imposer d'exigences
          vérifiables par un tiers~; par exemple COBIT, COSO
          (Ch.~\ref{chap:referentiels}).
    \item[Capacité d'un processus] Mesure de la manière dont un processus
          est mis en œuvre et fonctionne, de~0 à~5 dans COBIT~2019
          (Ch.~\ref{chap:cobit-capacite}).
    \item[Cascade des objectifs] Chaîne qui relie les besoins des parties
          prenantes, les objectifs d'entreprise, les objectifs d'alignement
          et les objectifs de gouvernance et de management
          (Ch.~\ref{chap:cobit-objectifs}).
    \item[Certification] Attestation, par un organisme tiers accrédité,
          qu'un système de management est conforme à une norme, sur un
          périmètre et pour une durée donnés (Ch.~\ref{chap:normes}).
    \item[Chaîne de déploiement] Suite automatisée des étapes par lesquelles
          un changement de code passe jusqu'à la mise en production
          (Ch.~\ref{chap:agile-cloud}).
    \item[Changement] Ajout, modification ou retrait de tout élément
          susceptible d'avoir un effet sur un service
          (Ch.~\ref{chap:itil-changements}).
    \item[Chronologie de défaillance] Suite datée des faits, décisions et
          absences de décision autour d'une défaillance, chaque entrée
          rattachée à une pièce (Ch.~\ref{chap:prologue}).
    \item[Co-création de valeur] Idée selon laquelle la valeur naît de la
          collaboration entre fournisseur et consommateur du service
          (Ch.~\ref{chap:itil-svs}).
    \item[Comité stratégique, comité de pilotage des SI] Le premier assiste
          le conseil sur les orientations, la valeur et les risques du
          numérique~; le second, au niveau de la direction générale,
          priorise les projets et arbitre les ressources
          (Ch.~\ref{chap:gouv-entreprise}).
    \item[Composante] Élément qui contribue au fonctionnement du système de
          gouvernance~; COBIT~2019 en distingue sept types
          (Ch.~\ref{chap:cobit-archi}).
    \item[Conformité] Respect d'exigences obligatoires (lois, règlements,
          contrats)~; elle se démontre (Ch.~\ref{chap:normes}).
    \item[Contrat de sous-traitance] Contrat entre le fournisseur et un tiers
          extérieur dont il dépend pour fournir le service
          (Ch.~\ref{chap:itil-niveaux}).
    \item[Contrôles d'application] Contrôles intégrés à un système, qui
          agissent sur le traitement des transactions (Ch.~\ref{chap:audit}).
    \item[Contrôles généraux informatiques] Contrôles qui s'appliquent à tout
          l'environnement informatique~: accès, changements, exploitation,
          développement et acquisition (Ch.~\ref{chap:gouv-entreprise},
          \ref{chap:audit}).
    \item[Coûts d'agence] Dépenses de surveillance du principal, dépenses
          d'engagement de l'agent, et perte résiduelle due aux divergences
          d'intérêts (Ch.~\ref{chap:gouv-entreprise}).
    \item[Coût total de possession (TCO)] Somme des coûts sur la durée de vie
          d'une solution, de l'acquisition à la sortie
          (Ch.~\ref{chap:alignement}).
    \item[Déclaration d'applicabilité] Liste des mesures de l'annexe~A de
          l'ISO/IEC~27001, retenues ou exclues, avec leur justification
          (Ch.~\ref{chap:normes}).
    \item[Défaillance de gouvernance] Absence, inadéquation ou contournement
          d'un mécanisme de décision, de contrôle ou de reddition de
          comptes, qui laisse un risque connu se réaliser
          (Ch.~\ref{chap:prologue}).
    \item[Demande de service] Demande d'un utilisateur portant sur une
          prestation déjà prévue dans le service (Ch.~\ref{chap:itil-incidents}).
    \item[DevOps] Pratiques qui rapprochent développement et exploitation,
          pour livrer des changements fréquents, fiables et réversibles
          (Ch.~\ref{chap:agile-cloud}).
    \item[Disponibilité] Aptitude d'un service à remplir sa fonction quand on
          en a besoin~; se mesure sur le temps de service convenu
          (Ch.~\ref{chap:itil-niveaux}).
    \item[Domaine d'intérêt] Thème de gouvernance traité par un ensemble
          d'objectifs et de composantes (Ch.~\ref{chap:cobit-archi}).
    \item[Droits de décision] Répartition formelle du pouvoir de décider sur
          une question donnée (Ch.~\ref{chap:prologue},
          \ref{chap:alignement}).
    \item[Durée maximale d'interruption admissible (DMIA)] Durée au-delà de
          laquelle l'interruption d'une activité devient inacceptable
          (Ch.~\ref{chap:itil-niveaux}).
    \item[Élément de configuration] Tout composant qui doit être géré pour
          fournir un service (Ch.~\ref{chap:itil-changements}).
    \item[Erreur connue] Problème analysé, dont la cause est identifiée, mais
          pas encore résolu (Ch.~\ref{chap:itil-incidents}).
    \item[Évaluation par les pairs] Examen d'un article par des chercheurs
          indépendants avant sa publication (Ch.~\ref{chap:recherche-mesure}).
    \item[Facteur de conception] Facteur qui influence la conception du
          système de gouvernance~; COBIT~2019 en retient onze
          (Ch.~\ref{chap:cobit-archi}).
    \item[Garantie] Assurance que le service répondra aux exigences
          convenues~: disponibilité, capacité, sécurité, continuité
          (Ch.~\ref{chap:itil-svs}).
    \item[Gestionnaire de processus] Personne qui pilote l'exécution
          quotidienne d'un processus (Ch.~\ref{chap:processus}).
    \item[Gouvernance d'entreprise] Relations, règles et mécanismes par
          lesquels une société est dirigée et contrôlée
          (Ch.~\ref{chap:gouv-entreprise}).
    \item[Gouvernance des données] Exercice de l'autorité et du contrôle sur
          la gestion des données (Ch.~\ref{chap:donnees-ia}).
    \item[Gouvernance des SI] Activités par lesquelles l'organe de
          gouvernance évalue, dirige et surveille l'usage de l'information
          et des technologies, pour qu'il crée de la valeur
          (Ch.~\ref{chap:prologue}, \ref{chap:gouv-entreprise}).
    \item[Habilitation des changements] Pratique qui maximise le nombre de
          changements réussis, en évaluant les risques, en autorisant et en
          maîtrisant le calendrier (Ch.~\ref{chap:itil-changements}).
    \item[Incident] Interruption non planifiée d'un service, ou dégradation
          de sa qualité (Ch.~\ref{chap:itil-incidents}).
    \item[Incident majeur] Incident dont l'impact exige une réponse
          dépassant le traitement ordinaire (Ch.~\ref{chap:itil-incidents}).
    \item[Indicateur clé d'objectif (KGI), de performance (KPI)] Le premier
          dit si un objectif est atteint, après coup~; le second, si un
          processus a des chances de l'atteindre, en cours de route
          (Ch.~\ref{chap:processus}).
    \item[Management des SI] Activités qui planifient, bâtissent, exploitent
          et surveillent les SI, selon la direction fixée par la gouvernance
          (Ch.~\ref{chap:prologue}).
    \item[Matrice RACI] Tableau qui croise activités et rôles~: qui réalise
          (R), qui en répond (A, un seul), qui est consulté (C), qui est
          informé (I) (Ch.~\ref{chap:processus}).
    \item[Maturité] Mesure de performance d'un domaine d'intérêt~: il atteint
          un niveau quand tous ses processus l'atteignent
          (Ch.~\ref{chap:cobit-capacite}).
    \item[Modèle des trois lignes] Répartition des rôles entre l'organe de
          gouvernance, le management et l'audit interne
          (Ch.~\ref{chap:audit}).
    \item[Niveau de capacité cible] Niveau de~0 à~5 que l'organisation décide
          d'atteindre pour un processus (Ch.~\ref{chap:cobit-conception}).
    \item[Norme certifiable] Norme qui formule des exigences dont le respect
          peut être audité par un organisme de certification
          (Ch.~\ref{chap:referentiels}).
    \item[Objectif d'alignement, d'entreprise] Le second est un but de
          l'organisation~; le premier, un but que l'informatique doit
          atteindre pour le soutenir (Ch.~\ref{chap:cobit-objectifs}).
    \item[Objectif de gouvernance ou de management] Ce qu'il faut atteindre
          pour que l'information et les technologies contribuent aux
          objectifs~; COBIT~2019 en compte quarante
          (Ch.~\ref{chap:cobit-archi}).
    \item[Optimisation des ressources, des risques] Disposer des ressources
          suffisantes et bien utilisées~; maintenir les risques dans
          l'appétence fixée par la gouvernance (Ch.~\ref{chap:alignement}).
    \item[Partie prenante] Groupe ou individu qui affecte les objectifs d'une
          organisation ou en est affecté (Ch.~\ref{chap:gouv-entreprise}).
    \item[Plan de retour arrière] Actions qui ramènent à l'état antérieur si
          un changement échoue, avec leurs critères et leur délai
          (Ch.~\ref{chap:itil-changements}).
    \item[Portefeuille] Ensemble des projets, programmes et services, géré
          comme un tout pour maximiser la valeur (Ch.~\ref{chap:alignement}).
    \item[Pratique (COBIT)] Action clé pour atteindre un objectif~; codée,
          par exemple DSS02.04 (Ch.~\ref{chap:cobit-objectifs}).
    \item[Prépublication] Version d'un travail rendue publique avant son
          évaluation par les pairs (Ch.~\ref{chap:recherche-mesure}).
    \item[Problème] Cause, réelle ou potentielle, d'un ou de plusieurs
          incidents (Ch.~\ref{chap:itil-incidents}).
    \item[Processus] Ensemble d'activités corrélées qui transforme des
          entrées en sorties, en vue d'un objectif (Ch.~\ref{chap:processus}).
    \item[Propriétaire de processus] Personne qui répond de la conception, de
          la performance et de l'amélioration d'un processus
          (Ch.~\ref{chap:processus}).
    \item[Ratio d'impact disproportionné] Taux de décisions favorables du
          groupe le moins favorisé divisé par celui du plus favorisé~; un
          ratio inférieur à~0,8 est un signal d'alerte
          (Ch.~\ref{chap:donnees-ia}).
    \item[Réalisation des bénéfices] Créer de la valeur grâce aux
          technologies, et éliminer ce qui n'en crée pas assez
          (Ch.~\ref{chap:alignement}).
    \item[Recueil de bonnes pratiques] Pratiques reconnues, présentées comme
          recommandations, que l'organisation adapte~; par exemple ITIL
          (Ch.~\ref{chap:referentiels}).
    \item[Redevabilité] Obligation de rendre compte d'une responsabilité
          reçue, et d'en assumer les conséquences (Ch.~\ref{chap:prologue}).
    \item[Relation d'agence] Contrat par lequel un principal délègue à un
          agent un pouvoir de décision (Ch.~\ref{chap:gouv-entreprise}).
    \item[RTO, RPO] Durée maximale de reprise d'un service~; perte de
          données maximale admissible, exprimée en durée
          (Ch.~\ref{chap:itil-niveaux}).
    \item[Service] Moyen de permettre la co-création de valeur, en
          facilitant les résultats que veut le client
          (Ch.~\ref{chap:itil-svs}).
    \item[Solution de contournement] Moyen de réduire l'effet d'un incident
          ou d'un problème sans en traiter la cause
          (Ch.~\ref{chap:itil-incidents}).
    \item[Source de première main] Document produit par une partie qui a eu
          accès direct aux faits et engage sa responsabilité sur leur
          exactitude (Ch.~\ref{chap:defaillances}).
    \item[Système de management] Ensemble d'éléments en interaction qui sert
          à établir des politiques et des objectifs, et à les atteindre
          (Ch.~\ref{chap:normes}).
    \item[Système de valeur des services] Selon ITIL~4, la manière dont les
          composantes et activités d'une organisation transforment une
          demande en valeur (Ch.~\ref{chap:itil-svs}).
    \item[Utilité] Ce que fait le service pour répondre à un besoin
          (Ch.~\ref{chap:itil-svs}).
    \item[Valeur] Bénéfices, utilité et importance perçus~; elle est perçue
          par celui qui la reçoit (Ch.~\ref{chap:itil-svs}).
    \item[Valeur actuelle nette (VAN)] Somme des flux d'un projet ramenés à
          aujourd'hui par un taux d'actualisation (Ch.~\ref{chap:alignement}).
    \item[Vérification des systèmes d'information] Selon le GUID~5100,
          examen des contrôles liés aux SI pour déceler des écarts à des
          critères (Ch.~\ref{chap:audit}).
\end{description}
